\section*{Bab \ref{ch:g2:from-seed-to-plant} --- Dari Biji Menjadi Tumbuhan}

\begin{solution}{exo:g2:from-seed-to-plant:1}
Sebuah \omterm{def:g1:plants-around-us:plant}{tumbuhan} mungil yang sudah dimulai (\omterm{def:g2:from-seed-to-plant:seed}{lembaga} \omterm{def:g1:plants-around-us:plant}{tumbuhannya}, dengan
bakal \omterm{def:g1:plants-around-us:parts}{akar} dan bakal tunas) dan cadangan makanan untuk memberinya
makan.
\end{solution}

\begin{solution}{exo:g2:from-seed-to-plant:2}
\omterm{def:g2:from-seed-to-plant:germination}{Berkecambah} berarti terbangun lalu mulai tumbuh: \omterm{def:g2:from-seed-to-plant:seed}{bijinya} membengkak,
\omterm{def:g1:the-five-senses:sense}{kulitnya} terbelah, dan ia mulai menjadi \omterm{def:g1:plants-around-us:plant}{tumbuhan}. \omterm{def:g1:plants-around-us:parts}{Akarnya} keluar lebih
dahulu, lalu menukik ke bawah.
\end{solution}

\begin{solution}{exo:g2:from-seed-to-plant:3}
Air dan kehangatan. Cahaya tidak ada dalam daftarnya --- \omterm{def:g2:from-seed-to-plant:seed}{bijinya}
\omterm{def:g2:from-seed-to-plant:germination}{berkecambah} dengan senang hati dalam gelap, di bawah tanah.
\end{solution}

\begin{solution}{exo:g2:from-seed-to-plant:4}
Ia hidup dari cadangan makanan yang dikemas di dalam \omterm{def:g2:from-seed-to-plant:seed}{bijinya} --- kedua
belahan tebal kacangnya. Cadangan itu memberi makan \omterm{prop:g2:from-seed-to-plant:cycle}{kecambahnya} sampai
\omterm{def:g1:plants-around-us:parts}{daun} hijaunya mencapai cahaya dan ia dapat makan sendiri.
\end{solution}

\begin{solution}{exo:g2:from-seed-to-plant:5}
Misalnya: kacang, ercis, miju-miju, beras, \omterm{def:g2:from-seed-to-plant:seed}{biji} \omterm{def:g1:plants-around-us:parts}{bunga} matahari, \omterm{def:g2:from-seed-to-plant:seed}{biji}
apel, \omterm{def:g2:from-seed-to-plant:seed}{biji} besar buah persik.
\end{solution}

\begin{solution}{exo:g2:from-seed-to-plant:6}
Karena sebutir \omterm{def:g2:from-seed-to-plant:seed}{biji} membutuhkan kehangatan selain air: di dalam tanah
musim dingin yang dingin ia akan menunggu, atau membusuk, alih-alih
\omterm{def:g2:from-seed-to-plant:germination}{berkecambah}. Musim semi membawa tanah hangat yang ditunggu \omterm{def:g2:from-seed-to-plant:seed}{bijinya}.
\end{solution}

\begin{solution}{exo:g2:from-seed-to-plant:7}
\omterm{def:g2:from-seed-to-plant:seed}{Biji}, \omterm{prop:g2:from-seed-to-plant:cycle}{kecambah}, tanaman berdaun lebat, \omterm{def:g1:plants-around-us:plant}{tumbuhan} berbunga, polong
berisi kacang baru.
\end{solution}

\begin{solution}{exo:g2:from-seed-to-plant:8}
Kedua \omterm{def:g1:plants-around-us:parts}{akarnya} tumbuh ke bawah dan kedua tunasnya tumbuh ke atas,
bagaimanapun kacangnya terbaring. \omterm{prop:g2:from-seed-to-plant:cycle}{Kecambahnya} selalu menemukan arah
atas dan bawah --- tepat seperti yang diramalkan catatannya.
\end{solution}

\begin{solution}{exo:g2:from-seed-to-plant:9}
Di dalam laci, \omterm{def:g2:from-seed-to-plant:seed}{biji-bijinya} tidak punya air, jadi tidak dapat
\omterm{def:g2:from-seed-to-plant:germination}{berkecambah}; tetapi \omterm{def:g2:from-seed-to-plant:seed}{biji} yang baik tanpa air tidak mati --- ia
menunggu. Setahun kemudian, ketika diberi air dan kehangatan, \omterm{def:g2:from-seed-to-plant:seed}{biji}
yang menunggu itu terbangun seperti biasa.
\end{solution}

\begin{solution}{exo:g2:from-seed-to-plant:10}
Cadangan makanan \omterm{def:g2:from-seed-to-plant:seed}{bijinya} --- kedua belahan tebal kacangnya ---
memainkan peran kuning telur. \omterm{def:g2:animals-grow-up:born}{Telur} dan \omterm{def:g2:from-seed-to-plant:seed}{biji} sama-sama permulaan
sebuah \omterm{def:g1:living-or-not:living}{makhluk hidup} baru, yang dikemas bersama makanan yang
dibutuhkan si muda sebelum ia dapat makan sendiri.
\end{solution}
